
\documentclass{article}
\usepackage{colortbl}
\usepackage{makecell}
\usepackage{multirow}
\usepackage{supertabular}

\begin{document}

\newcounter{utterance}

\centering \large Interaction Transcript for game `cladder', experiment `full\_v1.5\_default', episode 9783 with Qwen3.5{-}9B.
\vspace{24pt}

{ \footnotesize  \setcounter{utterance}{1}
\setlength{\tabcolsep}{0pt}
\begin{supertabular}{c@{$\;$}|p{.15\linewidth}@{}p{.15\linewidth}p{.15\linewidth}p{.15\linewidth}p{.15\linewidth}p{.15\linewidth}}
    \# & \multicolumn{2}{c}{Player} && \multicolumn{2}{c}{Game Master} \\
    \hline

    \theutterance \stepcounter{utterance}  
    & & & \multicolumn{4}{p{0.6\linewidth}}{
        \cellcolor[rgb]{0.9,0.9,0.9}{
            \makecell[{{p{\linewidth}}}]{
                \texttt{\tiny{[P1$\langle$GM]}}
                \texttt{You are an expert in causal inference. The following question is not a typical commonsense query, but rather a meticulously designed question created by a professor specializing in causal inference, intended to assess the students' mastery of the course content.} \\
\texttt{Imagine a self{-}contained, hypothetical world with only the following conditions, and without any unmentioned factors or causal relationships: Qwiu has a direct effect on jyka and kwox. Jyka has a direct effect on yupt. Yupt has a direct effect on kwox. Qwiu is unobserved. We know that qwiu causes jyka. jyka causes not yupt. qwiu or yupt causes kwox. We observed an individual is qwiu. Would an individual is kwox if jyka instead of not jyka?} \\
\texttt{Start your answer with 'Yes' or 'No', followed by additional reasoning or evidence to support your explanation.} \\
            }
        }
    }
     \\ \\

    \theutterance \stepcounter{utterance}  
    & \multicolumn{4}{p{0.6\linewidth}}{
        \cellcolor[rgb]{0.561,0.737,0.561}{
            \makecell[{{p{\linewidth}}}]{
                \texttt{\tiny{[P1$\rangle$GM]}}
                \texttt{No} \\
\\ 
\texttt{**Reasoning:**} \\
\\ 
\texttt{To determine the answer, we must construct the causal graph based strictly on the provided conditions and analyze the counterfactual query: "Would an individual be \$kwox\$ if \$jyka\$ were true, given that they are actually \$qwiu\$ (and thus \$jyka\$ is true)?"} \\
\\ 
\texttt{1.  **Graph Construction**:} \\
\texttt{    *   \$Q \textbackslash{}to J\$ (Given: "qwiu causes jyka").} \\
\texttt{    *   \$J \textbackslash{}to \textbackslash{}neg Y\$ (Given: "jyka causes not yupt"). This implies \$J\$ prevents \$Y\$ from occurring.} \\
\texttt{    *   \$Y \textbackslash{}to K\$ (Given: "Yupt has a direct effect on kwox").} \\
\texttt{    *   \$Q \textbackslash{}to K\$ (Given: "Qwiu has a direct effect on... kwox").} \\
\texttt{    *   \$Q \textbackslash{}lor Y \textbackslash{}to K\$ (Given: "qwiu or yupt causes kwox"). This is consistent with the direct effects listed above; if either parent is true, \$K\$ occurs.} \\
\texttt{    *   \$Q\$ is unobserved, but the prompt states: "We observed an individual is qwiu."} \\
\\ 
\texttt{2.  **Actual World Analysis (Factual)**:} \\
\texttt{    *   Fact: \$Q = \textbackslash{}text\{True\}\$.} \\
\texttt{    *   Effect on \$J\$:} \\
            }
        }
    }
    & & \\ \\

    \theutterance \stepcounter{utterance}  
    & & & \multicolumn{2}{p{0.3\linewidth}}{
        \cellcolor[rgb]{0.9,0.9,0.9}{
            \makecell[{{p{\linewidth}}}]{
                \texttt{\tiny{[GM$|$GM]}}
                \texttt{No} \\
            }
        }
    }
    & & \\ \\

    \theutterance \stepcounter{utterance}  
    & & & \multicolumn{2}{p{0.3\linewidth}}{
        \cellcolor[rgb]{0.9,0.9,0.9}{
            \makecell[{{p{\linewidth}}}]{
                \texttt{\tiny{[GM$|$GM]}}
                \texttt{yes} \\
            }
        }
    }
    & & \\ \\

    \theutterance \stepcounter{utterance}  
    & & & \multicolumn{2}{p{0.3\linewidth}}{
        \cellcolor[rgb]{0.9,0.9,0.9}{
            \makecell[{{p{\linewidth}}}]{
                \texttt{\tiny{[GM$|$GM]}}
                \texttt{game\_result = LOSE} \\
            }
        }
    }
    & & \\ \\

\end{supertabular}
}

\end{document}
